\documentclass[11pt]{amsart}

\usepackage[margin=1.15in]{geometry}
\usepackage{amsmath,amssymb,amsthm,mathtools}
\usepackage{microtype}
\usepackage{hyperref}
\usepackage[nameinlink,capitalize]{cleveref}

\hypersetup{
  colorlinks=true,
  linkcolor=blue,
  citecolor=blue,
  urlcolor=blue
}

\theoremstyle{plain}
\newtheorem{theorem}{Theorem}[section]
\newtheorem{proposition}[theorem]{Proposition}
\newtheorem{lemma}[theorem]{Lemma}
\newtheorem{corollary}[theorem]{Corollary}

\theoremstyle{definition}
\newtheorem{definition}[theorem]{Definition}
\newtheorem{hypothesis}[theorem]{Hypothesis}

\theoremstyle{remark}
\newtheorem{remark}[theorem]{Remark}

\numberwithin{equation}{section}

\title[Conditional equilibrium criterion]{A Conditional Equilibrium Criterion for Compact Ancient Navier--Stokes Dynamics}

\author{Guillem Duran Ballester}
\date{}

\subjclass[2020]{35Q30, 35B40, 35B41, 35B44, 37B25, 76D05}
\keywords{Navier--Stokes equations, ancient solutions, invariant measures, self-similar variables, Type I singularities, Liouville theorems}

\begin{document}

\begin{abstract}
We prove a conditional equilibrium criterion for compact ancient solutions of
the three-dimensional Navier--Stokes equations in backward self-similar
variables.  The assumption is formulated on a compact time-translation
invariant hull: there is a continuous Lyapunov functional whose dissipation is
nonnegative and vanishes precisely on stationary trajectories.  Under this
assumption, every invariant probability measure is supported on stationary
solutions.  If the compact hull is minimal and uniformly bounded in
\(L^3(\mathbb R^3)\), the stationary self-similar theorem of
Necas--Ruzicka--Sverak forces the hull to consist only of the zero trajectory.
We then record the resulting conditional assembly: a normalized collection has
no nonzero ancient solutions provided its minimal elements generate compact
nonzero hulls satisfying the Lyapunov hypothesis.  The construction of such a
Lyapunov functional is not asserted here; it is isolated as the remaining
analytic input.
\end{abstract}

\maketitle

\section{Introduction}

We consider the incompressible Navier--Stokes equations in backward
self-similar variables,
\begin{equation}\label{eq:centered}
   \partial_\tau V+(V\cdot\nabla)V+\nabla P
   =
   \Delta V-\frac12(V+y\cdot\nabla V),
   \qquad \operatorname{div}V=0,
\end{equation}
on \(\mathbb R^3\times\mathbb R\).  The equation is autonomous in
\(\tau\).  If \(V\) is a solution and \(s\in\mathbb R\), then
\[
   V^s(y,\tau)=V(y,\tau+s)
\]
is again a solution.

The compactness and invariant-measure construction for time translates is
recalled below.  That construction yields a stationary statistical identity,
but it does not imply that the underlying trajectories are stationary: the
barycenter of an invariant measure carries a quadratic covariance term.  The
additional input considered here is an equilibrium principle excluding
recurrent nonstationary compact motion.  This paper records a precise
conditional form of that principle.

The hypothesis used below is explicit.  We assume the existence of a
continuous Lyapunov functional on the compact invariant hull, with an exact
dissipation identity and with zero dissipation precisely at stationary
trajectories.  This hypothesis implies that invariant measures are supported
on equilibria.  Uniform \(L^3\) control then places stationary elements within
the scope of the theorem of Necas--Ruzicka--Sverak \cite{NRS1996}.  Thus the
paper is a conditional assembly theorem, not an unconditional proof of
equilibrium rigidity for all compact ancient Navier--Stokes dynamics.

\subsection{Role of the hypothesis}

The strict Lyapunov hypothesis is the only input in this paper not already
provided by the preceding compactness, minimality, and stationary-rigidity
arguments.  Compactness gives invariant measures, and stationary \(L^3\)
rigidity gives zero stationary profiles.  The remaining point is the
equilibrium-support assertion: an invariant measure on a compact nonlinear
Navier--Stokes hull must be supported on stationary trajectories.  The
functional \(E\) and dissipation \(D\) in \cref{hyp:lyapunov} encode precisely
that assertion.

\section{Compact ancient dynamics}

We use the following solution class, topology, and compact-hull conventions.
They are recalled here to fix notation for the present argument.  The locally
smooth topology is the standard Fr\'echet topology of uniform convergence of
all space-time derivatives on compact subsets.  It is induced by the
countable family of seminorms
\[
   \|U\|_{m,N}
   =
   \max_{|\alpha|+j\le m}
   \sup_{\substack{|y|\le N\\ |\tau|\le N}}
   |\partial_y^\alpha\partial_\tau^jU(y,\tau)|,
   \qquad m,N\in\mathbb N .
\]
The metric
\[
   d(U,V)=
   \sum_{m,N=1}^\infty
   2^{-(m+N)}
   \frac{\|U-V\|_{m,N}}{1+\|U-V\|_{m,N}}
\]
induces this topology.  Consequently each compact hull considered here is a
compact metric space.

All probability measures below are Borel probability measures on the relevant
compact metric hull \(K\), and \(C(K)\) denotes the real-valued continuous
functions on \(K\).  Weak convergence of probability measures is understood
against test functions in \(C(K)\).  Since \(K\) is compact and metric, every
Borel probability measure on \(K\) is Radon; in particular its closed support
has full measure.  A probability measure \(\mu\) is invariant if
\[
   \int_K F(\Theta_tW)\,d\mu(W)
   =
   \int_K F(W)\,d\mu(W)
   \qquad\hbox{for every }F\in C(K)\hbox{ and every }t\in\mathbb R .
\]
If \(B\subset K\) is Borel, the phrase ``\(\mu\) is supported on \(B\)''
means \(\mu(K\setminus B)=0\).  The closed support
\(\operatorname{supp}\mu\) is the smallest closed subset of \(K\) with full
\(\mu\)-measure, equivalently the set of points whose every open
neighborhood has positive \(\mu\)-measure.

After applying the whole-space Helmholtz--Leray projection \(\mathbb P\),
equation \eqref{eq:centered} becomes
\begin{equation}\label{eq:projected}
   \partial_\tau V
   =
   \Delta V-\frac12y\cdot\nabla V-\frac12V
   -\mathbb P\operatorname{div}(V\otimes V).
\end{equation}
Let \(S(a)\), \(a>0\), denote the semigroup generated by
\(\Delta-\frac12y\cdot\nabla-\frac12\).  The pressure reconstruction and the
passage between the projected mild equation and the local classical
Navier--Stokes equation are included in the local smoothing and pressure
reconstruction statement below.

The mild formulation is used only to specify the ancient solution class and
to make the time-translation action intrinsic.  The arguments below use the
locally smooth representatives in the compact hull and distributional
identities obtained from \eqref{eq:centered}.  The distribution
\(\mathbb P\operatorname{div}(V\otimes V)\) is interpreted by duality against
smooth test fields and as a tempered-distribution Fourier multiplier on each
time slice.  No boundedness of the Riesz transforms on \(L^\infty\) is used in
this paper.

\begin{definition}[Bounded mild ancient solution]\label{def:mild}
A vector field \(V:\mathbb R^3\times\mathbb R\to\mathbb R^3\) is a bounded
mild ancient solution of \eqref{eq:centered} if
\[
   V\in L^\infty(\mathbb R^3\times\mathbb R),\qquad \operatorname{div}V=0
\]
in distributions, and if for every \(\sigma<\tau\)
\begin{equation}\label{eq:mild}
   V(\tau)
   =
   S(\tau-\sigma)V(\sigma)
   -
   \int_\sigma^\tau
   S(\tau-s)\mathbb P\operatorname{div}(V\otimes V)(s)\,ds
\end{equation}
as an identity in the weak-* topology of \(L^\infty(\mathbb R^3)\).
\end{definition}

\begin{definition}[Compact invariant hull]\label{def:compact-hull}
Let \(K\) be a nonempty compact set of bounded mild ancient solutions, with
compactness understood in the locally smooth
\[
   C^\infty_{\mathrm{loc}}(\mathbb R^3\times\mathbb R)
\]
topology.  For \(t\in\mathbb R\), write
\[
   \Theta_tW=W^t,\qquad W^t(y,\tau)=W(y,\tau+t).
\]
The set \(K\) is invariant if \(\Theta_tK=K\) for every \(t\in\mathbb R\).  It
is minimal if it has no proper nonempty compact invariant subset.
\end{definition}

\begin{lemma}[Continuity of time translation]\label{lem:translation-continuity}
For each \(t\in\mathbb R\), the map \(\Theta_t:K\to K\) is a homeomorphism.
Moreover, for every \(W\in K\), the orbit map
\[
   \mathbb R\ni t\mapsto \Theta_tW\in K
\]
is continuous, and the action map
\[
   \mathbb R\times K\ni(t,W)\mapsto\Theta_tW\in K
\]
is continuous.
\end{lemma}

\begin{proof}
Continuity of each time-translation map follows directly from the definition
of the locally smooth topology.  Joint continuity follows because convergence
on compact cylinders is stable under bounded time shifts.  The identity
\(\Theta_t^{-1}=\Theta_{-t}\) gives the asserted homeomorphism property.
\end{proof}

\begin{definition}[Stationary trajectory]\label{def:stationary}
An element \(W\in K\) is stationary if \(\Theta_tW=W\) for every
\(t\in\mathbb R\).  For a time-independent trajectory \(W(y,\tau)=w(y)\), the
corresponding stationary self-similar profile equation is
\begin{equation}\label{eq:stationary}
   (w\cdot\nabla)w+\nabla p
   =
   \Delta w-\frac12(w+y\cdot\nabla w),
   \qquad \operatorname{div}w=0 .
\end{equation}
\end{definition}

\begin{lemma}[Distributions with zero gradient]\label{lem:zero-gradient-constant}
Let \(\Omega\subset\mathbb R^3\) be connected and open.  If
\(q\in\mathcal D'(\Omega)\) satisfies \(\nabla q=0\) in
\(\mathcal D'(\Omega)\), then \(q\) is constant in \(\Omega\), in the sense
that there exists \(c\in\mathbb R\) with \(q=c\) in \(\mathcal D'(\Omega)\).
\end{lemma}

\begin{proof}
Let \(\rho_\varepsilon\) be a Friedrichs mollifier.  If
\(\Omega'\Subset\Omega\) and
\(\varepsilon<\operatorname{dist}(\Omega',\partial\Omega)\), then
\[
   q_\varepsilon(x)
   =
   \langle q(y),\rho_\varepsilon(x-y)\rangle
\]
is well-defined and smooth for \(x\in\Omega'\).  Moreover,
\[
   \nabla q_\varepsilon
   =
   (\nabla q)*\rho_\varepsilon
   =
   0
   \qquad\hbox{on }\Omega' .
\]
Thus \(q_\varepsilon\) is constant on every connected component of
\(\Omega'\).  In particular, if \(B\Subset\Omega\) is a ball, then
\(q_\varepsilon\) is constant on \(B\) for all sufficiently small
\(\varepsilon\).  To see that the limit is also constant, let
\(\varphi\in C_c^\infty(B)\) satisfy \(\int_B\varphi\,dx=0\).  Since
\(q_\varepsilon\) is constant on \(B\),
\[
   \langle q_\varepsilon,\varphi\rangle=0 .
\]
Passing to the limit gives \(\langle q,\varphi\rangle=0\) for every such
\(\varphi\).  Choose \(\eta\in C_c^\infty(B)\) with \(\int_B\eta\,dx=1\), and
set \(c=\langle q,\eta\rangle\).  For arbitrary
\(\psi\in C_c^\infty(B)\), the test function
\(\psi-(\int_B\psi\,dx)\eta\) has integral zero; hence
\[
   \langle q,\psi\rangle
   =
   c\int_B\psi\,dx .
\]
Thus the restriction of \(q\) to \(B\) is the constant distribution \(c\).

If \(B_1,B_2\Subset\Omega\) are balls with \(B_1\cap B_2\ne\varnothing\), the
constants obtained on \(B_1\) and \(B_2\) agree, because both represent the
restriction of \(q\) to \(B_1\cap B_2\).  It remains only to justify that this
propagates through all of \(\Omega\).  Fix a ball \(B_0\Subset\Omega\), and
let \(U\) be the union of all balls \(B\Subset\Omega\) that can be joined to
\(B_0\) by a finite chain of compactly contained balls with consecutive
nonempty intersections.  The set \(U\) is open by construction.  It is also
closed relative to \(\Omega\): if \(x_j\in U\) and \(x_j\to x\in\Omega\),
choose a ball \(B_x\Subset\Omega\) containing \(x\); then \(x_j\in B_x\) for
large \(j\), so the chain reaching a ball containing \(x_j\), followed by
\(B_x\), reaches \(B_x\).  Since \(\Omega\) is connected and \(U\) is nonempty,
we have \(U=\Omega\).  Hence the local constants agree throughout \(\Omega\).
\end{proof}

\begin{lemma}[Stationary elements and the profile equation]
\label{lem:stationary-profile}
Let \(W\in K\).  Then \(W\) is stationary if and only if there is a smooth
divergence-free vector field \(w:\mathbb R^3\to\mathbb R^3\) such that
\[
   W(y,\tau)=w(y)
   \qquad\hbox{for all }(y,\tau)\in\mathbb R^3\times\mathbb R .
\]
For every such \(w\), there is a scalar distribution \(p\), unique modulo
additive constants, such that \(w\) solves \eqref{eq:stationary} in
\(\mathcal D'(\mathbb R^3)\).
\end{lemma}

\begin{proof}
If \(W\) is stationary, then \(\Theta_tW=W\) as elements of the locally smooth
function space for every \(t\in\mathbb R\).  Equality in this Hausdorff
function space is equality of the smooth representatives; hence
\[
   W(y,\tau+t)=W(y,\tau)
   \qquad\hbox{for all }y,\tau,t .
\]
Taking \(\tau=0\) and then writing \(t\) as an arbitrary time variable gives
\(W(y,t)=W(y,0)\).  Thus \(W\) is independent of time.  Conversely, if
\(W(y,\tau)=w(y)\), then \(W(y,\tau+t)=W(y,\tau)\) for every \(t\), so
\(\Theta_tW=W\) for every \(t\).

Set \(w(y)=W(y,0)\).  The local smoothness of \(W\) gives
\(w\in C^\infty(\mathbb R^3;\mathbb R^3)\), and the distributional condition
\(\operatorname{div}W=0\) gives \(\operatorname{div}w=0\).  By local smoothing
and pressure reconstruction, every bounded mild ancient solution in the compact
hull has, on each compact space-time cylinder, a smooth local pressure for
which \eqref{eq:centered} holds pointwise.  For each integer \(N\ge1\), apply
that result on the cylinder \(B_N\times(-1,1)\).  There is a smooth pressure
\(P_N(y,\tau)\) on this cylinder satisfying \eqref{eq:centered} with
\(V=W\).  Since the present
\(W\) is independent of \(\tau\), the vector field
\[
   \Delta w-\frac12(w+y\cdot\nabla w)-(w\cdot\nabla)w
\]
is independent of \(\tau\), and therefore \(\nabla_yP_N(y,\tau)\) is
independent of \(\tau\) on \(B_N\times(-1,1)\).  Setting
\(p_N(y)=P_N(y,0)\), the local identity reduces on \(B_N\) to
\[
   (w\cdot\nabla)w+\nabla p_N
   =
   \Delta w-\frac12(w+y\cdot\nabla w),
   \qquad \operatorname{div}w=0 .
\]
On overlaps the gradients of two such local pressures agree, because the
remaining terms in the displayed equation are global functions of \(w\).
Set \(c_1=0\).  For each integer \(R\ge2\), apply
\cref{lem:zero-gradient-constant} on
\(B_{R-1}\) to the distribution \(p_R-p_{R-1}-c_{R-1}\).  Its gradient
vanishes on \(B_{R-1}\), because both \(\nabla p_R\) and
\(\nabla p_{R-1}\) are equal there to
\[
   \Delta w-\frac12(w+y\cdot\nabla w)-(w\cdot\nabla)w .
\]
Therefore the distribution is constant on \(B_{R-1}\).
Choose \(c_R\) to be the negative of that constant.  Then
\[
   p_R+c_R=p_{R-1}+c_{R-1}
   \qquad\hbox{in }\mathcal D'(B_{R-1}) .
\]

Define the global distribution \(p\) as follows.  If
\(\varphi\in C_c^\infty(\mathbb R^3)\), choose an integer \(R\ge1\) such that
\(\operatorname{supp}\varphi\subset B_R\), and set
\[
   \langle p,\varphi\rangle
   =
   \langle p_R+c_R,\varphi\rangle .
\]
This definition is independent of the chosen sufficiently large \(R\), because
the preceding compatibility relation iterates: if \(S>R\) are integers, then
\[
   p_S+c_S=p_R+c_R
   \qquad\hbox{in }\mathcal D'(B_R).
\]
Indeed, one applies the preceding compatibility identity successively on
\(B_{S-1},B_{S-2},\ldots,B_R\), and then restricts the resulting distributional
identity to \(B_R\).
Hence \(p\) is a well-defined scalar distribution on \(\mathbb R^3\), and its
gradient agrees with the patched local pressure gradients.

Let \(\Phi\in C_c^\infty(\mathbb R^3;\mathbb R^3)\).  Choose an integer \(R\)
so large that \(\operatorname{supp}\Phi\subset B_R\).  Testing the local
identity on \(B_R\) against \(\Phi\), and using the definition of \(p\) on
\(B_R\), gives
\[
   \left\langle
      (w\cdot\nabla)w+\nabla p
      -\Delta w+\frac12(w+y\cdot\nabla w),
      \Phi
   \right\rangle=0 .
\]
Since this holds for every compactly supported smooth vector field \(\Phi\),
\eqref{eq:stationary} holds in \(\mathcal D'(\mathbb R^3)\).  If \(p\) and
\(\tilde p\) both have the required gradient, then
\(\nabla(p-\tilde p)=0\) in \(\mathcal D'(\mathbb R^3)\), so
\cref{lem:zero-gradient-constant} applied on the connected set
\(\mathbb R^3\) gives
\(p-\tilde p\) constant.
\end{proof}

\begin{theorem}[Invariant measures on compact hulls]\label{thm:invariant-measures}
Let \(K\) be a compact invariant hull and let \(W_0\in K\).  For \(T>0\), set
\[
   \mu_T=\frac1T\int_0^T\delta_{\Theta_sW_0}\,ds .
\]
Equivalently, \(\mu_T\) is the push-forward of normalized Lebesgue measure on
\([0,T]\) under the continuous orbit map \(s\mapsto\Theta_sW_0\); hence
\(\mu_T\) is a Borel probability measure on \(K\).
Every sequence \(T_n\to\infty\) has a subsequence for which \(\mu_{T_n}\)
converges weakly to an invariant Borel probability measure \(\mu\) on \(K\).
If \(K\) is minimal, then the support of every invariant Borel probability
measure on \(K\) is all of \(K\).
\end{theorem}

\begin{proof}
This is the Krylov--Bogolyubov averaging argument for a continuous flow on a
compact metric space.  Compactness gives subsequential weak limits of the orbit
averages; invariance follows by comparing the averages of a test function
\(\phi\) and \(\phi\circ\Theta_t\), whose difference is supported only on the
two time intervals of length \(|t|\) at the ends of \([0,T]\).  If \(K\) is
minimal, the support of an invariant measure is a nonempty closed invariant
subset, hence all of \(K\).
\end{proof}

\section{The equilibrium hypothesis}

\begin{hypothesis}[Strict Lyapunov functional]\label{hyp:lyapunov}
Let \(K\) be a compact invariant hull.  Assume there are continuous functions
\[
   E:K\to\mathbb R,\qquad D:K\to[0,\infty)
\]
such that, for every \(W\in K\) and every \(t\ge0\),
\begin{equation}\label{eq:lyapunov}
   E(\Theta_tW)-E(W)
   =
   -\int_0^t D(\Theta_sW)\,ds .
\end{equation}
The integral is the ordinary Lebesgue integral of the continuous function
\([0,t]\ni s\mapsto D(\Theta_sW)\).  Thus \eqref{eq:lyapunov} is an integral
identity along each orbit; no differentiability of \(E(\Theta_tW)\) in \(t\)
is assumed.
Assume also that
\begin{equation}\label{eq:zero-dissipation}
   D(W)=0
   \quad\hbox{if and only if}\quad
   W\hbox{ is stationary}.
\end{equation}
\end{hypothesis}

\begin{remark}
The paper proves the consequences of \cref{hyp:lyapunov}.  It does not claim
that such a functional has been constructed for the full Navier--Stokes
dynamics.  Constructing it, or replacing it by an equivalent rigidity
principle, is the remaining equilibrium problem.
\end{remark}

\begin{lemma}[Invariant measures annihilate dissipation]\label{lem:zero-diss}
Let \(K\) be a compact invariant hull satisfying \cref{hyp:lyapunov}.  If
\(\mu\) is an invariant Borel probability measure on \(K\), then
\[
   \int_K D(W)\,d\mu(W)=0 .
\]
\end{lemma}

\begin{proof}
Fix \(t>0\).  The functions \(E\) and \(D\) are continuous on compact \(K\),
so they are bounded and Borel measurable.  By
\cref{lem:translation-continuity}, \((s,W)\mapsto D(\Theta_sW)\) is continuous
on \([0,t]\times K\).  In particular it is Borel measurable and bounded on the
compact product \([0,t]\times K\).  With respect to the finite product measure
\(ds\otimes d\mu(W)\), all integrals below are finite, and Fubini's theorem
applies.

Integrating \eqref{eq:lyapunov} with respect to \(\mu\), and using invariance
of \(\mu\) with the continuous test function \(E\), gives
\[
   0
   =
   \int_K E(\Theta_tW)\,d\mu(W)-\int_KE(W)\,d\mu(W)
   =
   -\int_K\int_0^tD(\Theta_sW)\,ds\,d\mu(W).
\]
By Fubini's theorem and invariance again, now with the continuous test
function \(D\), for each fixed \(s\in[0,t]\),
\[
   \int_K\int_0^tD(\Theta_sW)\,ds\,d\mu(W)
   =
   \int_0^t\int_KD(\Theta_sW)\,d\mu(W)\,ds
   =
   t\int_KD(W)\,d\mu(W).
\]
Combining the two displayed identities gives
\[
   t\int_KD(W)\,d\mu(W)=0 .
\]
Since \(t>0\), the claim follows.
\end{proof}

\begin{lemma}[Zero dissipation gives stationary support]\label{lem:support-zero-d}
Let \(D:K\to[0,\infty)\) be continuous and let \(\mu\) be a Borel probability
measure on \(K\).  If
\[
   \int_KD(W)\,d\mu(W)=0,
\]
then \(\mu\) is supported on the closed set \(\{D=0\}\).
\end{lemma}

\begin{proof}
For each \(\varepsilon>0\), the set
\[
   K_\varepsilon=\{W\in K:\ D(W)\ge\varepsilon\}
\]
is closed, hence Borel.  Since \(D\ge\varepsilon\) on \(K_\varepsilon\),
\[
   0=\int_KD\,d\mu\ge\varepsilon\,\mu(K_\varepsilon).
\]
Thus \(\mu(K_\varepsilon)=0\) for every \(\varepsilon>0\).  Since \(D\) is
continuous, the zero set \(\{D=0\}\) is closed.  Its complement is the
countable union
\[
   K\setminus\{D=0\}
   =
   \bigcup_{n=1}^\infty K_{1/n}.
\]
Indeed, if \(D(W)>0\), then \(1/n\le D(W)\) for all sufficiently large \(n\).
Hence \(K\setminus\{D=0\}\) has \(\mu\)-measure zero.
\end{proof}

\begin{theorem}[Equilibrium support]\label{thm:equilibrium-support}
Let \(K\) be a compact invariant hull satisfying \cref{hyp:lyapunov}.  Then
every invariant Borel probability measure on \(K\) is supported on stationary
trajectories.
\end{theorem}

\begin{proof}
Let \(\mu\) be an invariant Borel probability measure on \(K\).  By
\cref{lem:zero-diss},
\[
   \int_KD(W)\,d\mu(W)=0 .
\]
By \cref{lem:support-zero-d}, \(\mu\) is supported on \(\{D=0\}\).  The zero
set of \(D\) is precisely the set of stationary trajectories by
\eqref{eq:zero-dissipation}.  Hence every invariant measure is supported on
stationary trajectories.
\end{proof}

\begin{corollary}[Minimal compact hulls are stationary]\label{cor:minimal-stationary}
Let \(K\) be a minimal compact invariant hull satisfying
\cref{hyp:lyapunov}.  Then every element of \(K\) is stationary.  Consequently
\(K\) consists of a single stationary trajectory.
\end{corollary}

\begin{proof}
Choose \(W_0\in K\) and a sequence \(T_n\to\infty\).  By
\cref{thm:invariant-measures}, after passing to a subsequence the averages
\[
   \mu_{T_n}
   =
   \frac1{T_n}\int_0^{T_n}\delta_{\Theta_sW_0}\,ds
\]
converge weakly to an invariant Borel probability measure \(\mu\) on \(K\).
Since \(K\) is minimal, the support assertion in
\cref{thm:invariant-measures} yields \(\operatorname{supp}\mu=K\).

By \cref{thm:equilibrium-support}, \(\mu\) is supported on stationary
trajectories.  By \cref{hyp:lyapunov}, the set of stationary trajectories is
\(\{D=0\}\), and this set is closed because \(D\) is continuous.  We claim
that \(\operatorname{supp}\mu\subset\{D=0\}\).  Indeed, if
\(W\in\operatorname{supp}\mu\setminus\{D=0\}\), then the open set
\[
   U=K\setminus\{D=0\}
\]
is an open neighborhood of \(W\).  By the definition of support,
\(\mu(U)>0\).  This contradicts \(\mu(K\setminus\{D=0\})=0\).  Therefore
\[
   K=\operatorname{supp}\mu
   \subset\{W\in K:\ W\hbox{ is stationary}\}.
\]
Hence every element of \(K\) is stationary.

If \(W\in K\) is stationary, then the singleton \(\{W\}\) is compact and
nonempty.  It is invariant because \(\Theta_tW=W\) for every \(t\).  Since
\(K\) has no proper nonempty compact invariant subset, minimality gives
\(K=\{W\}\).
\end{proof}

\section{Stationary self-similar rigidity}

\begin{theorem}[Stationary self-similar theorem]\label{thm:NRS}
Let \(w\in C^\infty(\mathbb R^3;\mathbb R^3)\cap L^3(\mathbb R^3)\) solve
\eqref{eq:stationary} in distributions for some scalar distribution \(p\),
with \(\operatorname{div}w=0\).  Then \(w\equiv0\).
\end{theorem}

\begin{proof}
This is the stationary self-similar Liouville theorem of
Necas--Ruzicka--Sverak \cite{NRS1996}.  Their theorem is stated for weak
\(L^3(\mathbb R^3)\) solutions of the stationary Leray self-similar system;
the smooth distributional formulation above is a special case.
\end{proof}

\begin{theorem}[Rigidity of compact tight minimal dynamics]\label{thm:no-compact-tight}
Let \(K\) be a minimal compact invariant hull satisfying
\cref{hyp:lyapunov}.  Assume that there is \(L<\infty\) such that
\[
   \sup_{W\in K}\sup_{\tau\in\mathbb R}
   \|W(\tau)\|_{L^3(\mathbb R^3)}\le L .
\]
Then every element of \(K\) is identically zero.
\end{theorem}

\begin{proof}
By \cref{cor:minimal-stationary}, \(K=\{W\}\) for a stationary trajectory.
By \cref{lem:stationary-profile}, there is a smooth divergence-free profile
\(w\) such that \(W(y,\tau)=w(y)\), and \(w\) solves \eqref{eq:stationary} in
the sense of distributions for some pressure.  Since \(W\in K\), the assumed
bound gives
\[
   \|w\|_{L^3(\mathbb R^3)}
   =
   \|W(\tau)\|_{L^3(\mathbb R^3)}
   \le L
   \qquad\hbox{for every }\tau\in\mathbb R .
\]
Thus \(w\in L^3(\mathbb R^3)\).  The hypotheses of \cref{thm:NRS} are
satisfied, and hence \(w\equiv0\).  Therefore \(W\) is the identically zero
trajectory, and every element of \(K\) is identically zero.
\end{proof}

\section{Conditional assembly}

The preceding theorem is a compact-dynamical rigidity statement.  We now
record the precise form in which it combines with the minimal-element
reduction.  The reduction itself is isolated as the hypothesis below so that
this paper only supplies the compact-equilibrium implication.

\begin{hypothesis}[Minimal compact reduction]\label{hyp:minimal-compact}
Let \(A\) be a collection of bounded mild ancient solutions of
\eqref{eq:centered}.  Assume the following.

\begin{enumerate}
\item If \(A\) is nonempty, then \(A\) contains a nonzero element \(V_*\) of
minimal \(L^\infty(\mathbb R^3\times\mathbb R)\) norm among all elements of
\(A\).

\item For every such minimal element \(V_*\), the closure of
\(\{V_*^s:\ s\in\mathbb R\}\) contains a minimal compact invariant hull \(K\)
which contains a nonzero trajectory.

\item Every such \(K\) satisfies \cref{hyp:lyapunov}.

\item For every such compact invariant hull \(K\), there is \(L_K<\infty\)
such that every \(W\in K\) satisfies
\[
   \sup_{\tau\in\mathbb R}\|W(\tau)\|_{L^3(\mathbb R^3)}\le L_K .
\]
\end{enumerate}
\end{hypothesis}

\begin{remark}
This hypothesis is formulated for normalized collections, such as the fixed
nonvanishing classes obtained in the minimal-element reduction.  In
particular, if \(A\) is nonempty, then the first item excludes the
identically zero trajectory from \(A\).
\end{remark}

\begin{theorem}[Conditional ancient Liouville theorem]\label{thm:conditional-liouville}
If \(A\) satisfies \cref{hyp:minimal-compact}, then \(A\) is empty.
\end{theorem}

\begin{proof}
Suppose that \(A\) is nonempty.  By the first part of
\cref{hyp:minimal-compact}, choose a nonzero minimal element \(V_*\in A\).
By the second part, the time-translation hull of \(V_*\) contains a minimal
compact invariant hull \(K\) containing a nonzero trajectory.  By the third
and fourth parts, \(K\) satisfies the hypotheses of
\cref{thm:no-compact-tight}.  Therefore every element of \(K\) is identically
zero.  This contradicts the second part of
\cref{hyp:minimal-compact}, which asserts that the same \(K\) contains a
nonzero trajectory.  The assumption that \(A\) is nonempty is therefore false.
\end{proof}

\begin{corollary}[Conditional Type I exclusion]\label{cor:type-I-exclusion}
Assume that a putative Type I singularity in a given class produces, through
the Type I blow-up procedure, a nonzero bounded mild ancient solution belonging
to a collection \(A\) satisfying
\cref{hyp:minimal-compact}.  Then no Type I singularity occurs in that class.
\end{corollary}

\begin{proof}
The nonzero centered ancient limit is the Type I blow-up output, and the local
concentration/nonvanishing input underlying the surrounding singularity
dichotomy is included in the hypotheses of the construction.  Thus the blow-up
construction would give a nonzero element of \(A\), while
\cref{thm:conditional-liouville} gives \(A=\emptyset\).  Since a set cannot
both be empty and contain the produced nonzero ancient solution, this
contradiction excludes the singularity.
\end{proof}

\section{Conclusion}

The conditional structure is as follows.  Compactness and invariant measures
reduce compact ancient dynamics to invariant probability measures.  A strict
Lyapunov identity forces all such measures to be supported on stationary
trajectories.  Uniform \(L^3\) control then invokes the stationary
self-similar theorem and eliminates every nonzero minimal compact hull.  Thus
the remaining analytic task is to prove the strict Lyapunov hypothesis, or an
equivalent equilibrium-support theorem, for compact hulls generated by minimal
ancient Navier--Stokes solutions.

\begin{thebibliography}{99}

\bibitem{NRS1996}
J.~Ne\v{c}as, M.~R\r{u}\v{z}i\v{c}ka, and V.~\v{S}ver\'ak,
\emph{On Leray's self-similar solutions of the Navier--Stokes equations},
Acta Math. \textbf{176} (1996), no.~2, 283--294,
doi:10.1007/BF02551584.

\end{thebibliography}

\end{document}
